% GENERATED FILE. Edit canonical lessons, then run npm run generate:books.
% canonical-source-hash: fnv1a64:900fe923c397d51f
% canonical-lessons: KA-C33-yocisu, KA-C33-artha-maadiko, KA-C33-oodu, KA-C33-bare

\chapter{Four Verbs of Mind and Page}
\label{ch:ka-mind-verbs}
\hlchaptermodality{33}

\begin{chapteropening}
\textbf{By the end of this chapter:} \emph{I can say I think, I understand, I read and I write in Kannada, name the suffix that turns a Sanskrit noun into a Kannada verb, use the ending that hands an action back to its doer, and give both of the front-of-the-word sound laws that separate Kannada from Tamil.}

Say \kn{ನಾನು} \kn{ಕನ್ನಡ} \kn{ಬರೆಯುತ್ತೇನೆ}, keep it apart from \kn{ಬರುತ್ತೇನೆ} one vowel away, explain the glide a vowel-final stem needs, take bare back to the Dravidian line-drawing root behind Tamil \ta{வரை}, set the v $\to$ b law beside the p $\to$ h law, and count this chapter’s four verbs onto Chapter 32’s six to reach hattu.
\end{chapteropening}

\section[yōcisu]{\kn{ಯೋಚಿಸು} (yōcisu) — \textquotedblleft{}to think,\textquotedblright{} a borrowed word on a Kannada frame}

\label{lesson:KA-C33-yocisu}



\begin{quote}
\emph{Nanage gottu} — \textquotedblleft{}known to me,\textquotedblright{} a word with no beads at all, so the knower sat outside it. Now the opposite case: a verb with all three beads, whose stem is not Kannada.
\end{quote}

\subsection*{You'll want to know: \kn{ಯೋಚಿಸು}}
\begin{quote}
\textbf{\kn{ಯೋಚಿಸು}} — \emph{yōcisu} — \textbf{to think}
\end{quote}

\begin{quote}
\textbf{\kn{ನಾನು} \kn{ಯೋಚಿಸುತ್ತೇನೆ}.} — \emph{nānu yōcisuttēne.} — \textquotedblleft{}I think.\textquotedblright{}
\end{quote}

The pronoun is optional. The noun beside the verb is \textbf{\kn{ಯೋಚನೆ}} (\emph{yōcane}), \textquotedblleft{}a thought.\textquotedblright{}

\begin{sounds}
\textbf{\kn{ಯೋ}} is \textbf{\kn{ಯ}} (\emph{ya}) with the long-\emph{ō} sign, then \textbf{\kn{ಚಿ}}, then \textbf{\kn{ಸು}} — three even beats, \emph{yō-ci-su}.
\end{sounds}

\begin{cousinweb}
Chapter 9 built \textbf{\kn{ಕ್ಷಮಿಸು}} (\emph{kṣamisu}, \textquotedblleft{}to forgive\textquotedblright{}) from the Sanskrit noun \textbf{\kn{ಕ್ಷಮಾ}} (\emph{kṣamā}) plus \textbf{‑\kn{ಇಸು}} (\emph{‑isu}), Kannada's verb-making suffix. \emph{Yōcisu} is the same machine: \textbf{\kn{ಯೋಚನೆ}} plus \emph{‑isu}.

And \emph{yōcane} is itself Sanskrit — \textbf{\dv{योजना}} (\emph{yojanā}), \textquotedblleft{}a joining, an arranging,\textquotedblright{} from the root \textbf{\dv{युज्}} (\emph{yuj}), \textquotedblleft{}\textbf{to yoke}\textquotedblright{}: the root behind \emph{yoga} and behind English \textbf{yoke}. Thinking, so told, is putting things together.
\end{cousinweb}

\begin{cousinweb}
Kannada's inherited think-word did not die. It changed jobs. \textbf{\kn{ನೆನೆ}} (\emph{nene}) is ordinary and alive, and means \textquotedblleft{}\textbf{to call to mind, to remember}.\textquotedblright{}

The Dravidian dictionaries reconstruct that root as \emph{\textbf{*nen-ay}}, \textquotedblleft{}\textbf{to think}\textquotedblright{} — and Tamil \textbf{\ta{நினை}} (\emph{ninai}) is the same word, doing Tamil's everyday thinking. Two sisters divided one root's work: Kannada thinks with a borrowed noun and remembers with the old verb; Tamil thinks with the old verb outright.
\end{cousinweb}

\begin{grammarlens}[title={Grammar Lens: a borrowed stem, the native machine}]
Nothing of \emph{yōcisu}'s origin shows in how it works:

\begin{quote}
\emph{yōcisu} + \emph{-utt-} + \emph{-ēne} $\to$ \textbf{\kn{ಯೋಚಿಸುತ್ತೇನೆ}}
\end{quote}

Stem, tense, person — the three slots that built \emph{iruttēne}, and the six person endings untouched: \emph{yōcisuttīya}, \emph{yōcisuttāne}, \emph{yōcisuttāḷe}, \emph{yōcisuttāre}. A word may come from anywhere; once \emph{‑isu} is on it, it is a Kannada verb.
\end{grammarlens}

\subsection*{Your turn}
Say these aloud:

\begin{itemize}
\raggedright
  \item \textquotedblleft{}yōcisu\textquotedblright{} — think — then \textquotedblleft{}yōcisuttēne\textquotedblright{}
  \item two builds, one suffix — \textquotedblleft{}kṣamā, isu … yōcane, isu\textquotedblright{}
  \item the old root's two lives — \textquotedblleft{}nene\textquotedblright{} … Tamil \textquotedblleft{}ninai\textquotedblright{}
  \item beadless, then three-bead — \textquotedblleft{}nanage gottu … yōcisuttēne\textquotedblright{}
\end{itemize}

\begin{tcolorbox}[breakable,colback=teal!4,colframe=teal!35!black,title={Before you move on}]
Say \textquotedblleft{}I think,\textquotedblright{} name the suffix that made it a verb, and name the Chapter 9 word built that way. (\emph{Yōcisuttēne}; \textbf{‑\kn{ಇಸು}}; \emph{kṣamisu}.) Which Sanskrit root sits inside \emph{yōcane}, and what does it mean? (\emph{\textbf{Yuj}} — \textquotedblleft{}to yoke,\textquotedblright{} the root of \emph{yoga}.) What does \kn{ನೆನೆ} mean, and what is its Tamil cousin doing? (\textbf{To remember}; \emph{ninai} is Tamil's everyday \textquotedblleft{}think.\textquotedblright{}) Why does \kn{ಗೊತ್ತು} take no \emph{‑ēne}? (\textbf{It is not a verb} — no person slot, so its knower goes into the dative.)
\end{tcolorbox}
\section[arthamāḍiko]{\kn{ಅರ್ಥಮಾಡಿಕೊ} (arthamāḍiko) — \textquotedblleft{}to understand,\textquotedblright{} in three welded pieces}

\label{lesson:KA-C33-artha-maadiko}



\begin{quote}
\emph{Yōcisuttēne} was one borrowed noun with one Kannada suffix. The next verb is three separate words fused into one, and only the first is borrowed.
\end{quote}

\subsection*{You'll want to know: \kn{ಅರ್ಥಮಾಡಿಕೊ}}
\begin{quote}
\textbf{\kn{ಅರ್ಥಮಾಡಿಕೊ}} — \emph{arthamāḍiko} — \textbf{to understand}
\end{quote}

\begin{quote}
\textbf{\kn{ನಾನು} \kn{ಅರ್ಥಮಾಡಿಕೊಳ್ಳುತ್ತೇನೆ}.} — \emph{nānu arthamāḍikoḷḷuttēne.} — \textquotedblleft{}I understand.\textquotedblright{}
\end{quote}

Long, and worth saying slowly, because every piece has a name:

\begin{itemize}
\raggedright
  \item \textbf{\kn{ಅರ್ಥ}} (\emph{artha}) — \textquotedblleft{}\textbf{meaning},\textquotedblright{} a Sanskrit noun taken in whole.
  \item \textbf{\kn{ಮಾಡಿ}} (\emph{māḍi}) — \textquotedblleft{}\textbf{having done},\textquotedblright{} from the native \textbf{\kn{ಮಾಡು}} (\emph{māḍu}, \textquotedblleft{}to do\textquotedblright{}) of Chapter 5's \emph{kelasa māḍu}, \textquotedblleft{}to work.\textquotedblright{}
  \item \textbf{‑\kn{ಕೊ}} (\emph{‑ko}) — \textquotedblleft{}\textbf{and take it for yourself}.\textquotedblright{}
\end{itemize}

Meaning — made — taken back. That is the whole word.

\begin{sounds}
\textbf{\kn{ರ್ಥ}} stacks a silent \emph{r} over \textbf{\kn{ಥ}} (\emph{tha}, breathy), and the long form holds its \textbf{\kn{ಳ್ಳ}} — \emph{koḷ-ḷut-tē-ne}, the retroflex \emph{ḷ} doubled.
\end{sounds}

\begin{grammarlens}[title={Grammar Lens: ‑\kn{ಕೊ}, the ending that gives the action back}]
\textbf{‑\kn{ಕೊ}} is a worn-down \textbf{\kn{ಕೊಳ್ಳು}} (\emph{koḷḷu}), a real verb meaning \textquotedblleft{}to take \textbf{for oneself}.\textquotedblright{} Glued onto another verb it says: whatever this action produced, the doer keeps. Tamil \textbf{\ta{கொள்}} (\emph{koḷ}) and Telugu \textbf{\te{కొను}} (\emph{konu}) are the same ancient word at the same job.

Kannada has a second way to say it, with no doer at all:

\begin{quote}
\textbf{\kn{ನನಗೆ} \kn{ಅರ್ಥವಾಯಿತು}.} — \emph{nanage arthavāyitu.} — \textquotedblleft{}\textbf{to me} it became meaning\textquotedblright{} — \textquotedblleft{}\textbf{I understood}.\textquotedblright{}
\end{quote}

That is Chapter 6's dative subject, the frame that carried \emph{nanage gottu}. So Kannada offers a choice: take the meaning yourself with \emph{‑ko}, or let it arrive at you in the dative.
\end{grammarlens}

\begin{cousinweb}
Telugu builds understanding from the same parts: \textbf{\te{అర్థం} \te{చేసుకో}} (\emph{arthaṃ cēsuko}) — \emph{arthaṃ}, its own \textquotedblleft{}do\textquotedblright{} verb, its own \emph{‑ko}. Tamil declines the loan for a native verb, \textbf{\ta{புரிந்துகொள்}} (\emph{purintukoḷ}), yet keeps the \emph{koḷ} — the same instinct that left Tamil's everyday \textquotedblleft{}think\textquotedblright{} as the inherited \emph{ninai} while Kannada reached for a Sanskrit noun.
\end{cousinweb}

\subsection*{Your turn}
Say these aloud:

\begin{itemize}
\raggedright
  \item the three pieces — \textquotedblleft{}artha … māḍi … ko\textquotedblright{}
  \item \textquotedblleft{}arthamāḍikoḷḷuttēne\textquotedblright{} — I understand
  \item the dative way — \textquotedblleft{}nanage arthavāyitu\textquotedblright{}
  \item two datives together — \textquotedblleft{}nanage gottu … nanage arthavāyitu\textquotedblright{}
  \item this chapter so far — \textquotedblleft{}yōcisuttēne … arthamāḍikoḷḷuttēne\textquotedblright{}
\end{itemize}

\begin{tcolorbox}[breakable,colback=teal!4,colframe=teal!35!black,title={Before you move on}]
Name the three pieces of \kn{ಅರ್ಥಮಾಡಿಕೊ} and say which one is borrowed. (\emph{Artha}, \emph{māḍi}, \emph{‑ko}; \textbf{artha} is Sanskrit, the other two are native.) What does \emph{‑ko} do, and which verb is it worn down from? (\textbf{Hands the action back to its doer}; \emph{koḷḷu}.) Give the other way to say \textquotedblleft{}I understood,\textquotedblright{} and name the Chapter 6 sentence it matches. (\emph{Nanage arthavāyitu}; \emph{nanage gottu}.) Which sister keeps the ending but refuses the noun? (\textbf{Tamil} — \emph{purintukoḷ}.)
\end{tcolorbox}
\section[ōdu]{\kn{ಓದು} (ōdu) — \textquotedblleft{}to read,\textquotedblright{} and a warning about one letter}

\label{lesson:KA-C33-oodu}



\begin{quote}
\emph{Arthamāḍikoḷḷuttēne} — four syllables of ending. The next verb is two letters long, and the whole difficulty is keeping it apart from one you already have.
\end{quote}

\subsection*{You'll want to know: \kn{ಓದು}}
\begin{quote}
\textbf{\kn{ಓದು}} — \emph{ōdu} — \textbf{to read}, and also \textbf{to study}
\end{quote}

\begin{quote}
\textbf{\kn{ನಾನು} \kn{ಕನ್ನಡ} \kn{ಓದುತ್ತೇನೆ}.} — \emph{nānu kannaḍa ōduttēne.} — \textquotedblleft{}I read Kannada.\textquotedblright{}
\end{quote}

Say it beside Chapter 32's \textbf{\kn{ನೋಡು}} (\emph{nōḍu}, \textquotedblleft{}to see\textquotedblright{}). They differ by a single opening consonant and by which \emph{d} they use:

\noindent
\begin{tabularx}{\linewidth}{@{}>{\raggedright\arraybackslash}X>{\raggedright\arraybackslash}X>{\raggedright\arraybackslash}X@{}}
\toprule
\textbf{Kannada} & \textbf{said} & \textbf{means} \\
\midrule
\textbf{\kn{ಓದು}} & \emph{ōdu}, plain \emph{d} & read \\
\textbf{\kn{ನೋಡು}} & \emph{nōḍu}, curled \emph{ḍ} & see \\
\bottomrule
\end{tabularx}

\emph{Ōdu} also does the work English splits off as \textquotedblleft{}study.\textquotedblright{} Kannada does not keep the two apart: the same verb covers reading a page and reading for a degree.

\begin{sounds}
\textbf{\kn{ಓ}} is the independent long \emph{ō}, a letter in its own right rather than a sign hung on a consonant. Then \textbf{\kn{ದು}} (\emph{du}) — tongue at the teeth, not curled back. That is the one contrast to guard: \emph{nōḍu} curls, \emph{ōdu} does not.
\end{sounds}

\begin{cousinweb}
\emph{Ōdu} is inherited, not borrowed. The Dravidian dictionaries reconstruct \emph{\textbf{*ōtu}}, \textquotedblleft{}\textbf{to recite, to read},\textquotedblright{} and list the sisters: Tamil \textbf{\ta{ஓது}} (\emph{ōtu}) and Malayalam \textbf{\ml{ഓതുക}} (\emph{ōtuka}).

Notice the meaning that comes first in that gloss. The old sense is \textbf{reciting aloud} — reading in a world where you always said the words out loud — and Tamil's word stayed there, narrowing to the chanting of sacred verse. Tamil's everyday \textquotedblleft{}read\textquotedblright{} became a different word, \textbf{\ta{படி}} (\emph{paṭi}); Telugu's is \textbf{\te{చదువు}} (\emph{caduvu}).

So Kannada is the sister that left the inherited word in its ordinary job. That is the reverse of what happened with thinking, where Kannada took the loan and Tamil kept the old root as \emph{ninai}. Neither language borrows on principle; each word made its own bargain.
\end{cousinweb}

\subsection*{Your turn}
Say these aloud:

\begin{itemize}
\raggedright
  \item the pair, slowly — \textquotedblleft{}ōdu … nōḍu\textquotedblright{}
  \item \textquotedblleft{}nānu kannaḍa ōduttēne\textquotedblright{} — I read Kannada
  \item the sisters' read-words — \textquotedblleft{}ōdu … ōtu … paṭi … caduvu\textquotedblright{}
  \item the chapter so far — \textquotedblleft{}yōcisuttēne, arthamāḍikoḷḷuttēne, ōduttēne\textquotedblright{}
\end{itemize}

\begin{tcolorbox}[breakable,colback=teal!4,colframe=teal!35!black,title={Before you move on}]
Say \textquotedblleft{}I read Kannada,\textquotedblright{} then say \textquotedblleft{}I see\textquotedblright{} — and name what separates the two verbs. (\emph{Nānu kannaḍa ōduttēne}; \emph{nōḍuttēne}; the opening \emph{n‑} and a \textbf{curled \emph{ḍ}} against a \textbf{plain \emph{d}}.) What did the old root mean before it meant reading? (\textbf{Reciting aloud} — and Tamil's \emph{ōtu} stayed there.) Which sister keeps the inherited word for everyday reading? (\textbf{Kannada}.) And which one keeps the inherited word for everyday thinking? (\textbf{Tamil}, with \emph{ninai}.)
\end{tcolorbox}
\section[bare]{\kn{ಬರೆ} (bare) — \textquotedblleft{}to write,\textquotedblright{} which began as scratching a line}

\label{lesson:KA-C33-bare}



\begin{quote}
\emph{Ōduttēne}, \textquotedblleft{}I read.\textquotedblright{} Its partner sits one vowel from a verb you have had since Chapter 32.
\end{quote}

\subsection*{You'll want to know: \kn{ಬರೆ}}
\begin{quote}
\textbf{\kn{ಬರೆ}} — \emph{bare} — \textbf{to write}, equally \textbf{to draw}
\end{quote}

\begin{quote}
\textbf{\kn{ನಾನು} \kn{ಕನ್ನಡ} \kn{ಬರೆಯುತ್ತೇನೆ}.} — \emph{nānu kannaḍa bareyuttēne.} — \textquotedblleft{}I write Kannada.\textquotedblright{}
\end{quote}

The trap: \textbf{\kn{ಬರು}} (\emph{baru}) is the come-stem under \emph{bā}.

\noindent
\begin{tabularx}{\linewidth}{@{}>{\raggedright\arraybackslash}X>{\raggedright\arraybackslash}X@{}}
\toprule
\textbf{Kannada} & \textbf{means} \\
\midrule
\textbf{\kn{ಬರೆಯುತ್ತೇನೆ}} \emph{bareyuttēne} & I write \\
\textbf{\kn{ಬರುತ್ತೇನೆ}} \emph{baruttēne} & I come \\
\bottomrule
\end{tabularx}

Chapter 20's weather line takes the second — \textbf{\kn{ಮಳೆ} \kn{ಬರುತ್ತಿದೆ}} (\emph{maḷe baruttide}), "rain is coming.

\begin{grammarlens}[title={Grammar Lens: the glide a vowel-final stem needs}]
Why \emph{bare‑\textbf{y}‑uttēne}? The stem ends in a vowel, and Kannada will not let two vowels collide, so it slips a \textbf{\kn{ಯ}} (\emph{y}) between:

\begin{quote}
\emph{bare} + \textbf{y} + \emph{-utt-} + \emph{-ēne} $\to$ \textbf{\kn{ಬರೆಯುತ್ತೇನೆ}}
\end{quote}

Stems in \emph{‑u} need none: \emph{ōdu} gives \emph{ōduttēne} straight off. Stems in \emph{‑e} or \emph{‑i} take the glide. The six person endings are unchanged: \emph{bareyuttīri}, \emph{bareyuttāne}, \emph{bareyuttāre}.
\end{grammarlens}

\begin{cousinweb}
\emph{Bare} is inherited, from a root reconstructed as \emph{\textbf{*warV-}}, \textquotedblleft{}\textbf{to scratch, to draw lines}.\textquotedblright{} Every sister still means \textbf{draw} as readily as \textbf{write} — Tamil \textbf{\ta{வரை}} (\emph{varai}), Malayalam \textbf{\ml{വരയ്ക്കുക}} (\emph{varaykkuka}), Telugu \textbf{\te{రాయు}} (\emph{rāyu}).

Now the front of those words: Tamil \emph{v‑}, Kannada \emph{b‑} — Chapter 32's law, the one that gave \textbf{\kn{ಬಾ}} (\emph{bā}) for \emph{vā}. Kannada has two such laws:

\noindent
\begin{tabularx}{\linewidth}{@{}>{\raggedright\arraybackslash}X>{\raggedright\arraybackslash}X>{\raggedright\arraybackslash}X@{}}
\toprule
\textbf{law} & \textbf{Tamil or older} & \textbf{Kannada} \\
\midrule
\emph{p} $\to$ \emph{h} & \emph{pattu}, \emph{pasir} & \emph{hattu}, \emph{hasiru} \\
\emph{v} $\to$ \emph{b} & \emph{vā}, \emph{varai} & \emph{bā}, \emph{bare} \\
\bottomrule
\end{tabularx}

Between them they explain much of Kannada's distinctive sound.
\end{cousinweb}

\begin{grammarlens}[title={What you've built: four more verbs, and ten in all}]
Chapter 32 handed you \textbf{\kn{ಆರು}} (\emph{āru}) verbs: \emph{iru}, \emph{hōgu}, \emph{baru}, \emph{tinnu}, \emph{nōḍu}, \emph{gottu}. Four more makes \textbf{\kn{ಹತ್ತು}} (\emph{hattu}), ten, and none needed a new machine. What they added was history: \emph{yōcisu} a Sanskrit noun on a Kannada suffix, \emph{arthamāḍiko} one welded to two native verbs, \emph{ōdu} and \emph{bare} inherited outright.
\end{grammarlens}

\subsection*{Your turn}
Say these aloud:

\begin{itemize}
\raggedright
  \item \textquotedblleft{}nānu kannaḍa bareyuttēne\textquotedblright{} — I write Kannada
  \item the pair that must not merge — \textquotedblleft{}bareyuttēne … baruttēne … maḷe baruttide\textquotedblright{}
  \item the root across the sisters — \textquotedblleft{}bare … varai … rāyu\textquotedblright{}
  \item the two laws — \textquotedblleft{}pattu, hattu; pasir, hasiru; vā, bā; varai, bare\textquotedblright{}
  \item the four, then the count — \textquotedblleft{}yōcisuttēne, arthamāḍikoḷḷuttēne, ōduttēne, bareyuttēne … āru, hattu\textquotedblright{}
\end{itemize}

\begin{tcolorbox}[breakable,colback=teal!4,colframe=teal!35!black,title={Before you move on}]
Say \textquotedblleft{}I write Kannada,\textquotedblright{} then \textquotedblleft{}I come.\textquotedblright{} (\emph{Nānu kannaḍa bareyuttēne}; \emph{baruttēne} — \textbf{one vowel} apart.) Why the \emph{y}? (\textbf{The stem ends in a vowel}; \emph{ōdu} needs none.) What did \emph{*warV‑} mean, and what is Tamil's form? (\textbf{To scratch a line}; \emph{varai}, its \emph{v} against Kannada's \emph{b}.) Run the chapter's four and count them onto Chapter 32's six. (\emph{Yōcisuttēne, arthamāḍikoḷḷuttēne, ōduttēne, bareyuttēne} — \emph{āru} plus four is \emph{hattu}.)
\end{tcolorbox}
